\subsection{波动方程不等于完整的运动答案}
质点运动必须给初位置、初速度；连续介质则要对每一个位置给初位移、初速度。波动方程只规定变化规律，初始条件才决定是哪一次运动，边界条件又决定传播到端点后的反射方式。
对无界弦，右行项 $f(x-ct)$ 在 $x$ 增加 $c\Delta t$、时间增加 $\Delta t$ 后保持相同的自变量，所以波形向右；左行项 $g(x+ct)$ 同理向左。固定材料点只随时间上下振动，不能把波形的水平平移理解成弦本身向右滑动。

\begin{origexamplebox}{松手后分裂成两份的弦脉冲}\label{ch12:ex:teaching-initial-pulse}
一根足够长的均匀弦，波速为 $c$。在 $t=0$ 时拉成位移形状 $F(x)$ 并从静止释放，即 $y(x,0)=F(x)$、$y_t(x,0)=0$；$F$ 为局域的平滑函数。求后续波形，并解释为什么两列波的振幅各为原来的一半，能量却没有丢失。
\end{origexamplebox}
\begin{xesolution}
写 $y=f(x-ct)+g(x+ct)$。初位移给 $f(x)+g(x)=F(x)$；初速度给 $-cf'(x)+cg'(x)=0$。后式说明 $f-g$ 为常量。局域条件使远处位移趋零，常量可取零，故
\[
 y(x,t)=\frac12F(x-ct)+\frac12F(x+ct).
\]
每一份以速度 $c$ 向相反方向传播。代回 $t=0$，两半重叠恢复 $F$，两者的速度贡献反号，确实使初速度为零。
设弦张力为 $\mathcal T$、线密度为 $\mu$，满足 $\mathcal T=\mu c^2$。初始全是弹性势能：
\[
 E_0=\frac{\mathcal T}{2}\int_{-\infty}^{\infty}[F'(x)]^2\dd x.
\]
对右行半幅脉冲，$y_x=F'/2$、$y_t=-cF'/2$，动能与势能相等，其总能量为
\[
 E_+=\int\left(\frac{\mu c^2}{8}[F']^2+
 \frac{\mathcal T}{8}[F']^2\right)\dd x=\frac12E_0.
\]
左行脉冲同样带走 $E_0/2$。把“一半振幅对应四分之一能量”直接套在初始静止弦形上，会漏掉行波中的动能；初始状态与单向行波的动能条件不同。该解还要求端点反射尚未返回所研究区域。
\end{xesolution}
